\chapter{Conclusiones y trabajos futuros}\label{cap:conclusiones}

En este capítulo se resumen las conclusiones del trabajo, organizadas
según los objetivos planteados en la Sección~\ref{sec:objetivos}, y se
proponen las líneas de trabajo futuro que se derivan de los resultados y
de las limitaciones identificadas en el Capítulo~\ref{cap:resultados}.

\section{Conclusiones}\label{sec:conclusiones}
Este trabajo ha abordado la navegación reactiva sin mapa de un AGV en una
rejilla bidimensional con obstáculos estáticos mediante aprendizaje por
refuerzo profundo. Se ha formulado el problema como un POMDP, se ha
construido un entorno de simulación con un LiDAR simulado y una función
de recompensa compuesta, y se han entrenado y evaluado agentes DQN y PPO
sobre cuatro mapas de dificultad creciente. A lo largo de diecisiete
configuraciones experimentales, seis de ellas con cinco semillas de
entrenamiento, se han identificado los modos de fallo de la política
reactiva y se ha aplicado a cada uno un remedio específico. Las
conclusiones principales son las siguientes.

La observabilidad parcial no es solo un marco formal, sino la causa de
un modo de fallo concreto. La formulación como POMDP
(Sección~\ref{sec:pomdp}) anticipaba que una observación local no
identifica el estado. Los experimentos lo confirman: una política
determinista sin memoria queda atrapada en ciclos entre celdas con
observaciones casi idénticas, y esos ciclos explican la diferencia entre
evaluar la política de forma voraz o muestreando sus acciones. Añadir a
la observación el último desplazamiento del robot, una memoria de un
solo paso, reduce los ciclos de $42{,}6$ a $6{,}2$ por cada $1\,000$
episodios y eleva el éxito \textit{greedy} de \texttt{PPOAgent} del
$95{,}7\%$ al $99{,}4\%$, con cinco semillas. Con esta memoria las dos
modalidades de evaluación coinciden, de modo que la política determinista
puede desplegarse sin perder rendimiento.

Una observación local compacta es suficiente. La observación de
$18$ componentes, ocho sectores del LiDAR más la distancia y el rumbo a
la meta, permite aprender rutas prácticamente óptimas: la mediana de los
pasos de los episodios con éxito coincide con la del camino mínimo.
Triplicar el alcance del sensor no cambia el rendimiento, por lo que la
percepción no limita el problema estudiado. Una memoria mínima elegida
con criterio resulta, además, más eficaz que una memoria genérica: la
política recurrente con GRU no eliminó los ciclos y quedó por debajo del
agente sin memoria.

La función de recompensa de diseño propio es adecuada, y sus variantes no
la mejoran. La recompensa compuesta de progreso, coste por paso,
penalización por proximidad y términos terminales permite aprender en
los cuatro mapas. Ni el progreso geodésico ni el progreso con signo
mejoraron el resultado, y el segundo refutó la hipótesis de que los
ciclos se debieran a la forma de la recompensa. El currículo por etapas
aceleró la fase inicial del aprendizaje sin cambiar el rendimiento final.

El enmascarado de acciones inválidas elimina las colisiones, y con ellas
el principal límite del mapa más difícil. El techo de
\texttt{simple\_map}, en torno al $78\%$, se debía a choques junto a
paredes finas y no a mínimos locales, como se había supuesto
inicialmente. Con el enmascarado, ese mapa supera el $98\%$ de éxito en
la modalidad \textit{sampled}. Este resultado debe leerse con cautela: la
máscara usa la ocupación real de las celdas vecinas y actúa como un
escudo externo, no como una capacidad aprendida ni como una garantía
formal de seguridad.

PPO supera a DQN de forma robusta en este problema. Con la misma
observación, el mismo enmascarado y la misma red, \texttt{PPOAgent}
alcanza un $99{,}4 \pm 0{,}4\%$ de éxito \textit{greedy} frente al
$87{,}9 \pm 3{,}8\%$ de Double DQN. La ventaja se mantiene con el PPO
estándar, con las dos modalidades de evaluación y con o sin memoria
implícita. DQN muestra además sobreestimación de los valores $Q$ al
comienzo del entrenamiento, rutas menos eficientes y una dispersión entre
semillas unas diez veces mayor que la de PPO.

La implementación de PPO condiciona el resultado. El PPO
estándar aprende antes que la variante con \textit{buffer}, pero converge
prematuramente si no se limita la divergencia entre políticas sucesivas.
Con un KL objetivo de $0{,}015$ se queda a unos dos puntos de
\texttt{PPOAgent}, que obtiene el mejor resultado del trabajo pese a
apartarse del algoritmo original. Por ello, la conclusión principal, PPO
frente a DQN, se ha comprobado con las dos implementaciones.

La metodología experimental es parte del resultado. Varias
conclusiones intermedias cambiaron al revisar los experimentos: el fallo
de \texttt{simple\_map} no era un mínimo local, el orden entre las
variantes de DQN de una sola semilla era ruido y el PPO de referencia no
era el algoritmo del artículo original. Repetir las configuraciones
finales con cinco semillas, evaluar todos los agentes sobre los mismos
episodios y distinguir entre la variabilidad de la evaluación y la del
entrenamiento ha sido necesario para que las comparaciones sean fiables.

Estas conclusiones se circunscriben a un robot de cuatro acciones
absolutas con cinemática discreta, a cuatro mapas estáticos que se usan
tanto en el entrenamiento como en la evaluación, a los hiperparámetros
documentados y a un presupuesto de $40\,000$ episodios por ejecución. Los
hiperparámetros no se ajustaron de forma sistemática, lo que afecta
especialmente a DQN, del que solo se probaron dos tasas de aprendizaje.

\section{Trabajos futuros}\label{sec:trabajos_futuros}
Los resultados y las limitaciones anteriores sugieren las siguientes
líneas de continuación, ordenadas de la más inmediata a la más ambiciosa.

La primera línea es un currículo con una transición de dificultad
gradual entre mapas. El conjunto actual de mapas presenta un salto brusco de dificultad: de
tres mapas con obstáculos grandes y compactos se pasa a
\texttt{simple\_map}, con obstáculos finos, pasillos estrechos y
esquinas. El currículo probado en este trabajo ordenaba esos mismos
cuatro mapas y la distancia a la meta, y por eso solo podía acelerar el
comienzo del aprendizaje. La propuesta consiste en construir una familia de
mapas cuya dificultad crezca de forma controlada, por ejemplo mediante
generación procedimental, y medir esa dificultad con indicadores
objetivos: la densidad de ocupación, la anchura mínima de los pasillos,
el grosor de los obstáculos, el número de giros del camino mínimo o el
cociente entre la distancia geodésica y la euclídea. El entrenamiento
avanzaría por esa escala cuando el agente superase un umbral de éxito,
y podría combinarse con la selección automática de las posiciones de
partida en función del valor estimado por el crítico, como propone
\textit{NavACL} \citep{xue2022using}. Cabe esperar que una transición
gradual reduzca el número de episodios necesarios en los mapas
difíciles y, sobre todo, que mejore la generalización.

Ese currículo debería acompañarse de una evaluación de la generalización
en mapas no vistos. Todos los agentes se han evaluado en los mismos mapas en los que se entrenaron, con
posiciones de partida y meta nuevas. El paso siguiente es reservar un
conjunto de mapas de prueba que la política no haya visto, de dificultad
comparable a la de los mapas de entrenamiento, para medir si la política
reactiva ha aprendido a navegar o ha memorizado la geometría de los
escenarios. Esta evaluación es el complemento natural del currículo
anterior.

Otra línea es obtener la seguridad sin información privilegiada. El
enmascarado actual consulta la ocupación real de las celdas vecinas. Una versión desplegable
debería construir la máscara a partir de las lecturas del propio sensor o
de un sensor de seguridad de corto alcance, con el ruido que ello
implica, o bien sustituirla por mecanismos de aprendizaje por refuerzo
seguro con restricciones explícitas. Comparar la política entrenada con y
sin máscara bajo un escudo construido así indicaría cuánta capacidad de
evitar obstáculos aprende realmente la red.

También conviene explorar una memoria más rica y mejor entrenada. La
memoria de un paso basta en estos mapas, pero podría no bastar en escenarios con ambigüedades más
largas, como pasillos simétricos. Conviene estudiar historias de varios
desplazamientos en la observación y repetir el estudio de la política
recurrente con enmascarado, más semillas y estrategias de inicialización
del estado oculto que reproduzcan mejor el estado de la red al actuar.

Los hiperparámetros y las variantes de los algoritmos merecen un ajuste
sistemático. Un barrido de hiperparámetros con el mismo presupuesto para
cada algoritmo haría más sólida la comparación entre DQN y PPO, en
especial en la tasa de aprendizaje de DQN. También quedan por evaluar la
normalización de las observaciones, variantes de DQN como
\textit{Dueling DQN} o la repetición priorizada de experiencia, y las
variantes de PPO ya preparadas en el código: decaimiento lineal de la
tasa de aprendizaje, recorte dual, recorte con retroceso, penalización
adaptativa de la divergencia KL e inicialización reducida de la última
capa de la política.

El paso siguiente en cuanto al vehículo es adoptar cinemáticas más
realistas. Todos los experimentos usan el robot de cuatro acciones
absolutas. El entorno ya incluye un robot con orientación
y acciones de avance y giro, que no se ha entrenado, y la
Sección~\ref{sec:cinematica_agv} describe modelos continuos de
tracción diferencial, de Ackermann y de AGV híbrido triciclo-diferencial
\citep{sanchez2022modelado}. Una progresión natural es entrenar primero
con orientación en la rejilla, después con cinemática continua y, por
último, con la dinámica del vehículo y condiciones de adherencia
variables.

Por último, quedan los escenarios más complejos y la validación en un
robot real. Quedan fuera del alcance de este trabajo los obstáculos dinámicos, la
coordinación entre varios AGV y la transferencia de la simulación a un
vehículo físico. Para esta última, un paso intermedio sería un simulador
físico con sensores realistas, antes de la validación en un AGV
industrial como el estudiado por \citet{sanchez2022modelado}.
